% Generated draft: requires mathematical and source review.
\documentclass[11pt]{article}
\usepackage[margin=25mm]{geometry}
\usepackage{fontspec}
\setmainfont{Latin Modern Roman}
\usepackage{amsmath,amssymb,mathtools,unicode-math}
\setmathfont{Latin Modern Math}
\usepackage{enumitem,needspace,xurl,hyperref}
\setlength{\parindent}{0pt}
\setlength{\parskip}{5pt}
\setlength{\emergencystretch}{3em}
\newcommand{\N}{\mathbb{N}}
\newcommand{\Z}{\mathbb{Z}}
\newcommand{\Q}{\mathbb{Q}}
\newcommand{\R}{\mathbb{R}}
\newcommand{\C}{\mathbb{C}}
\newcommand{\F}{\mathbb{F}}
\newcommand{\A}{\mathbb{A}}
\newcommand{\PP}{\mathbb{P}}
\newcommand{\E}{\mathbb{E}}
\newcommand{\eps}{\varepsilon}
\newcommand{\dd}{\,\mathrm{d}}
\begin{document}
\section*{UnsolvedMath ID: 6700038}
\textbf{Source problem code:} catalogue entry
\subsection{Definitions and mathematical statement}
Let \(n\geq5\) and \(X=\operatorname{SO}(n)=\{A\in M_n(\R):A^{\mathsf T}A=I,\ \det A=1\}\). Let \(\bar g\) be a bi-invariant Riemannian metric on \(X\), meaning invariant under both left and right translations. The bar distinguishes this reference metric from a competing smooth metric \(g\); the extraction prints both as \(g\), which would otherwise make the strict scalar curvature comparison impossible.

With the usual quadratic-form order on metrics, the question is whether there exists \(g\) such that
\[
 g_x(v,v)\geq\bar g_x(v,v)\quad(x\in X,\ v\in T_xX),\qquad
 \operatorname{Scal}_g(x)>\operatorname{Scal}_{\bar g}(x)\quad(x\in X).
\]
Equivalently, can one rule out this simultaneous metric enlargement and everywhere strict scalar curvature improvement? The source says only that no known method rules it out, in its 2017 discussion; that is neither a proof of existence nor a verified statement of present knowledge. The competing metric is not assumed left invariant or bi-invariant.

The bi-invariant reference comes from the supplied summary and the adjacent local record ID 6700037. That context discusses area extremality, whereas the present extracted inequality is \(g\geq\bar g\). The displayed formulation uses the usual tensor inequality; whether an area comparison was intended instead needs source checking. The reference metric's normalization is not supplied.
\subsection{Short English statement}
For \(\operatorname{SO}(n)\), \(n\geq5\), can a metric dominate a bi-invariant reference metric while having strictly larger scalar curvature everywhere? The source describes a gap in available methods rather than asserting that such a metric exists.
\subsection{Sources}
\begin{itemize}
\item Local UnsolvedMath registry, permanent ID 6700038, source code catalogue entry; Misha Gromov, 2017; update 2026-08-06. Adjacent local ID 6700037, source code AMR-066-0037, supplies the bi-invariant-metric and area-extremality context.
\item Gromov, \emph{101 Questions, Problems and Conjectures around Scalar Curvature} (2017), printed tag [?41], occurrence 40, PDF page 48: \url{https://www.ihes.fr/~gromov/wp-content/uploads/2018/08/101-problemsOct1-2017.pdf}. Registry access date 2026-07-29; normalized PDF extraction from an incomplete public author-hosted PDF; current status and reuse terms need review.
\item The supplied summary reports the case as open and connects the difficulty to trivial tangent bundles and Dirac-index methods. It supplies no additional exact reference or verified present-day resolution.
\end{itemize}
\end{document}
